\begin{figure}[htbp]
\centering
\begin{tikzpicture}[
  catbox/.style={
    rectangle, rounded corners=3pt,
    minimum width=4.2cm, minimum height=1.1cm,
    draw=#1!60, fill=#1!8, line width=0.8pt,
    font=\bfseries\small, align=center
  },
  exbox/.style={
    rectangle, rounded corners=2pt,
    minimum width=4.2cm,
    draw=#1!30, fill=white, line width=0.5pt,
    font=\scriptsize, align=left, text=black,
    inner sep=5pt
  },
  detectbox/.style={
    rectangle, rounded corners=2pt,
    minimum width=4.2cm, minimum height=0.7cm,
    draw=#1!40, fill=#1!5, line width=0.5pt,
    font=\scriptsize\itshape, align=center, text=black
  },
]

% Column 1: Syntaktisch
\node[catbox=blue] (syn) {Syntaktisch};
\node[exbox=blue, below=4mm of syn] (syn-ex) {%
  \textbullet~Fehlendes \texttt{alt}-Attribut\\
  \textbullet~\texttt{<table>} ohne \texttt{<th>}\\
  \textbullet~Fehlender \texttt{aria-label}%
};
\node[detectbox=blue, below=3mm of syn-ex] (syn-det) {Regelbasiert pr\"ufbar};

% Column 2: Semantisch
\node[catbox=orange, right=8mm of syn] (sem) {Semantisch};
\node[exbox=orange, below=4mm of sem] (sem-ex) {%
  \textbullet~\texttt{alt="Bild"}\\
  \textbullet~Label: \textit{Eingabe}\\
  \textbullet~Linktext: \textit{hier klicken}%
};
\node[detectbox=orange, below=3mm of sem-ex] (sem-det) {Sprachverst\"andnis n\"otig};

% Column 3: Layout
\node[catbox=teal, right=8mm of sem] (lay) {Layout};
\node[exbox=teal, below=4mm of lay] (lay-ex) {%
  \textbullet~Unzureichender Kontrast\\
  \textbullet~Fehlender Fokusindikator\\
  \textbullet~Eingeschr\"ankte Zoomf\"ahigkeit%
};
\node[detectbox=teal, below=3mm of lay-ex] (lay-det) {Teils algorithmisch pr\"ufbar};

% Header
\node[above=8mm of sem, font=\small\bfseries] (title) {Taxonomie der Barrierefreiheitsverst\"o\ss e};

% Detection difficulty arrow
\draw[-{Stealth[length=3mm]}, line width=1pt, black]
  ([yshift=-3mm]syn-det.south west) -- ([yshift=-3mm]lay-det.south east)
  node[midway, below=1mm, font=\scriptsize, text=black] {Zunehmende Komplexit\"at der Erkennung};

\end{tikzpicture}
\caption{Taxonomie der Barrierefreiheitsverst\"o\ss e nach Fathallah et al. mit Beispielen und Erkennungscharakteristik je Kategorie}
\label{fig:taxonomie_verstoesse}
\end{figure}
